\section{The Executable Skill}\label{sec:skill}
% Principal structure:
%
% 4.1 From method description to executable procedure — a skill is a
%     versioned instruction+tool bundle the agent loads and OPERATES;
%     the method's rules become the agent's operating rules (log your own
%     activities; claims at most ai-confirmed; never fabricate telemetry).
% 4.2 ORCID-first setup — the only question asked of the human is their
%     ORCID iD; identity (name, affiliation) is resolved from the public
%     registry, not guessed; base IRI derives from the git remote; all
%     later commands auto-detect it. Figure~\ref{fig:workflow}.
% 4.3 CI/CD as the deterministic Build agent — every push: validate the
%     graph (invariants enforced server-side), render the dashboard,
%     compile the paper to PDF; the current paper version is always
%     available as a build artifact.
% 4.4 Versioning the method itself — the skill tree is the repository;
%     released snapshots are immutable (archive/), changes logged
%     (CHANGELOG), so the meta-experiment can cite the exact skill
%     version in force during each activity.

\begin{figure}[t]
\centering
% TODO(F3): TikZ workflow loop: ORCID question -> init (seed graph,
% register human, scaffold CI) -> register AI agent -> work loop
% (log / claim / source --verify) -> validate -> push -> CI (validate,
% dashboard, PDF artifact) -> audit/promote (human rungs).
\fbox{\parbox[c][4.2cm][c]{0.9\linewidth}{\centering\itshape
F3 placeholder: skill workflow from ORCID question to CI artifact}}
\caption{The skill's operating loop from first setup question to
continuously published paper build.}
\label{fig:workflow}
\end{figure}
